\DocumentMetadata{
  pdfversion=2.0,
  pdfstandard=ua-2,
  lang=en-US,
  tagging=on,
  tagging-setup={math/setup=mathml-SE,tabsorder=structure}
}
\documentclass[11pt,letterpaper]{article}
\usepackage[margin=0.68in,includefoot]{geometry}
\usepackage{newcomputermodern}
\usepackage{amsmath,amscd,mathtools}
\usepackage{tikz,adjustbox}
\providecommand{\square}{\mdlgwhtsquare}
\usepackage[hidelinks]{hyperref}
\hypersetup{
  pdftitle={Combinatorics PhD Exam, August 2017},
  pdfauthor={University of Florida Department of Mathematics},
  pdfsubject={Graduate mathematics examination},
  pdfdisplaydoctitle=true,
  bookmarksopen=true
}
\setcounter{secnumdepth}{0}
\setlength{\parindent}{0pt}
\setlength{\parskip}{0.28em}
\setlength{\emergencystretch}{3em}
\setlength{\leftmargini}{2.15em}
\setlength{\leftmarginii}{2.65em}
\setlength{\leftmarginiii}{3em}
\pagestyle{empty}
\raggedbottom
\allowdisplaybreaks
\NewTaggingSocketPlug{block/list/label}{examproblem}{%
  \tagstructbegin{tag=Lbl}\tagstructend%
  \tagstructbegin{tag=\LBody}%
  \tagstructbegin{tag=\ProblemHeadingTag,title-o={\ProblemHeadingText}}%
  \tagmcbegin{tag=\ProblemHeadingTag,actualtext={\ProblemHeadingText}}%
  #2%
  \tagmcend%
  \tagstructend%
}
\newenvironment{examproblems}[2]{%
  \def\ProblemHeadingTag{#2}%
  \AssignTaggingSocketPlug{block/list/label}{examproblem}%
  \begin{enumerate}%
  \setcounter{enumi}{#1}%
  \renewcommand{\labelenumi}{\textbf{\theenumi.}}%
  \setlength{\topsep}{0.35em}%
  \setlength{\partopsep}{0pt}%
  \setlength{\itemsep}{0.48em}%
  \setlength{\parsep}{0.18em}%
}{\end{enumerate}}
\newenvironment{examsubparts}{%
  \AssignTaggingSocketPlug{block/list/label}{default}%
  \begin{enumerate}%
  \setlength{\topsep}{0.18em}%
  \setlength{\partopsep}{0pt}%
  \setlength{\itemsep}{0.16em}%
  \setlength{\parsep}{0.1em}%
}{\end{enumerate}}
\newenvironment{examinstructions}{%
  \par\begingroup\itshape
}{%
  \par\endgroup\vspace{0.18em}
}
\makeatletter
\renewcommand\subsection{\@startsection{subsection}{2}{0pt}%
  {-1.25ex plus -.25ex minus -.1ex}%
  {0.55ex plus .1ex}%
  {\normalfont\large\bfseries}}
\renewcommand{\@maketitle}{%
  \newpage\null\vskip -1em%
  {\footnotesize\bfseries
    \href{https://www.ufl.edu/}{University of Florida}, \href{https://math.ufl.edu/}{Department of Mathematics}\par}%
  \vskip -0.5em%
  \tagstructbegin{tag=H1,title={Combinatorics PhD Exam, August 2017}}%
  \tagpdfparaOff%
  \tagmcbegin{tag=H1}%
    {\LARGE\bfseries\href{https://gma.math.ufl.edu/exams/combinatorics-phd/}{Combinatorics PhD Exam}, August 2017\par}%
  \tagmcend%
  \tagpdfparaOn%
  \tagstructend%
  \vskip 0.25em%
  \tagmcbegin{artifact}%
    \hrule height 1pt%
  \tagmcend%
  \vskip 1em%
}
\makeatother
\title{Combinatorics PhD Exam, August 2017}
\author{}
\date{}
\begin{document}
\maketitle
\thispagestyle{empty}
\begin{examproblems}{0}{H2}
\def\ProblemHeadingText{Problem 1.}
\item
Let \(a_n\) denote the number of strings of length~\(n\) of the letters A, B, C, and D such that the letter A appears an odd number of times.
\begin{examsubparts}
\item[\textnormal{(a)}]
Find a closed formula for \(a_n\).
\item[\textnormal{(b)}]
Find the exponential generating function for the sequence \(\{a_n\}\). (using no summation signs)
\end{examsubparts}
\def\ProblemHeadingText{Problem 2.}
\item
The set \([n] \times [n]\) is partially ordered by the relation \((a,b) \preceq (c,d)\) which holds when \(a \leq c\) and \(b \leq d\). Find, with proof, the length of a maximum chain and the length of a maximum antichain.

\par
What is the size of a minimum chain decomposition of this poset. (Recall that a chain decomposition of a poset is a partition of its elements into disjoint chains.)
\def\ProblemHeadingText{Problem 3.}
\item
This problem has two parts.
\begin{examsubparts}
\item[\textnormal{(a)}]
Prove that, for a simple graph~\(G\) with at least 5 vertices, at least one of~\(G\) or its complement has a cycle.
\item[\textnormal{(b)}]
Suppose that you color the edges of \(K_n\) using 2 colors. Show that there exists a monochromatic spanning tree.
\end{examsubparts}
\def\ProblemHeadingText{Problem 4.}
\item
A parity check matrix~\(H\) of a binary linear code~\(C\) can be defined as a generator matrix of the dual code \(C^\perp\).

\par
Show that~\(C\) is the null space of the transpose \(H^T\), where multiplication by \(H^T\) is on the right, i.e., \(cH^T\).

\par
Show that the minimum distance of the code equals the cardinality of a minimum dependent set of columns of~\(H\).

\par
If a generator matrix of a binary linear code~\(C\) is
\[
\begin{pmatrix}
1 & 0 & 1 & 0 & 1\\
0 & 1 & 1 & 1 & 0
\end{pmatrix},
\]
 then show that \((c_1,c_2,c_3,c_4,c_5)\) is a codeword if and only if (modulo 2)
\[
\begin{aligned}
c_1+c_2+c_3&=0\\
c_2+c_4&=0\\
c_1+c_5&=0.
\end{aligned}
\]
 What is the minimum distance of this code?
\def\ProblemHeadingText{Problem 5.}
\item
Let \(g(n)\) be the number of permutations of length~\(n\) in which each cycle is of even length. Find the exponential generating function of the sequence \(g(n)\), where \(n=0,1,2,\ldots\).
\def\ProblemHeadingText{Problem 6.}
\item
Let \(t_n\) be the total number of cycles of all permutations of length~\(n\). So \(t_1=1\), \(t_2=3\), and \(t_3=11\). Find an explicit formula for the numbers \(t_n\). Your answer can contain one summation sign.
\def\ProblemHeadingText{Problem 7.}
\item
Prove that the language \(\{a^{n^2}:n\in\mathbb{N}\}\) is not regular.
\def\ProblemHeadingText{Problem 8.}
\item
Let~\(\beta\) be a permutation of length~\(k\). Prove that there are precisely \(k^2+1\) permutations of length \(k+1\) containing~\(\beta\).
\end{examproblems}
\end{document}
